\documentclass[11pt,a4paper]{article}

\usepackage{kotex}
\usepackage{fontspec}
\setmainfont{Times New Roman}
\setmainhangulfont{AppleMyungjo}
\setsanshangulfont{Apple SD Gothic Neo}

\usepackage[a4paper,top=18mm,bottom=18mm,left=20mm,right=20mm]{geometry}
\usepackage{fancyhdr}
\setlength{\headheight}{24pt}
\usepackage{enumitem}
\usepackage{mdframed}
\usepackage{array}
\usepackage{booktabs}

\setlength{\parindent}{1em}
\linespread{1.18}

\pagestyle{fancy}
\fancyhf{}
\renewcommand{\headrulewidth}{0.6pt}
\fancyhead[L]{\sffamily\bfseries 2026 고1 영어 변형 --- 지문 36 정답 및 해설}
\fancyhead[R]{\sffamily 빙하기와 농업의 시작}
\renewcommand{\footrulewidth}{0pt}
\fancyfoot[C]{\framebox[16mm][c]{\thepage}}

\newcommand{\qno}[1]{\par\vspace{10pt}\noindent{\sffamily\bfseries #1.}\ }
\newmdenv[linewidth=0.6pt,roundcorner=3pt,innertopmargin=7pt,innerbottommargin=7pt,
          innerleftmargin=9pt,innerrightmargin=9pt,skipabove=5pt,skipbelow=5pt]{pbox}

\begin{document}

\begin{center}
{\fontsize{15}{17}\selectfont\sffamily\bfseries [지문 36] 정답 및 해설}
\end{center}
\vspace{8pt}

% ===== 정답 표 =====
\begin{center}
\renewcommand{\arraystretch}{1.4}
\begin{tabular}{cccc}
\toprule
\textbf{문항} & \textbf{유형} & \textbf{정답} & \textbf{배점} \\
\midrule
1 & 글 순서 배열 & ② & 2점 \\
2 & 인과 추론 & ① & 2점 \\
3 & 서술형 --- 표현 찾기 & (모범답안 참조) & 2점 \\
4 & 영어 제목 & ① & 2점 \\
\bottomrule
\end{tabular}
\end{center}

\vspace{10pt}

% ===== Q1 해설 =====
\qno{1}\textbf{[글 순서 배열]} \textbf{정답: ② (B) -- (A) -- (C)}

\vspace{4pt}
\noindent\textbf{정답 근거 및 흐름 분석:}

\begin{itemize}[leftmargin=2em,itemsep=3pt,topsep=3pt]
  \item \textbf{Opening:} ``During the Ice Age $\ldots$ it was much drier.'' --- 빙하기의 두 가지 특징(추위 + 건조함)을 제시한다.
  \item \textbf{(B):} ``In Ireland $\ldots$ there is far less evaporation, fewer clouds, and less rain.'' --- 개막 단락의 ``much drier''를 구체적으로 설명한다. 추위가 왜 건조함으로 이어지는지 인과 관계를 제시하는 연결 단락이다.
  \item \textbf{(A):} ``In this type of climate, farming isn't an option $\ldots$'' --- ``this type of climate''는 (B)에서 설명한 춥고 건조한 기후를 가리킨다. 이어 ``As the temperature rose $\ldots$''로 기후 변화의 전환점을 제시한다.
  \item \textbf{(C):} ``The world got warmer and wetter $\ldots$'' --- (A) 끝의 기온 상승과 빙하 해빙을 이어 받아, 더 따뜻하고 습한 세계에서 사람들이 농업을 점진적으로 시작했음을 설명한다.
\end{itemize}

\vspace{4pt}
\noindent\textbf{오답 소거:}
\begin{itemize}[leftmargin=2em,itemsep=2pt]
  \item ① (A)--(C)--(B): (A)의 ``In this type of climate''가 맥락 없이 첫 번째로 오면 지시 대상이 불분명하다.
  \item ③ (B)--(C)--(A): (C)의 따뜻하고 습한 세계 이후에 (A)의 ``In this type of climate''가 나오면, 그 표현이 춥고 건조한 기후를 가리키기 어려워진다.
  \item ④ (C)--(A)--(B): (C)가 첫 번째로 오면 ``warmer and wetter''의 배경 설명이 없어 맥락이 부족하다.
  \item ⑤ (C)--(B)--(A): (C)와 (B)의 순서가 뒤바뀌면 인과 관계가 역전된다.
\end{itemize}

\bigskip

% ===== Q2 해설 =====
\qno{2}\textbf{[인과 추론]} \textbf{정답: ①}

\vspace{4pt}
\noindent\textbf{정답 근거:} (B) 단락은 추위가 단순히 ``습함''으로 이어진다는 통념을 반박한다. 실제로 매우 추우면 증발이 줄고, 그 결과 구름과 비도 줄어든다. 이어 식물이 자라기 어려운 춥고 건조한 환경이라는 결론이 나온다. 따라서 ``매우 낮은 온도 → 대기 중 수분 이동 감소 → 구름 형성과 강수 감소''의 흐름을 나타낸 ①이 정답이다.

\vspace{4pt}
\noindent\textbf{오답 소거:}
\begin{itemize}[leftmargin=2em,itemsep=2pt]
  \item ② 낮은 온도가 Ireland를 더 습하게 만든다는 내용은 (B)의 반박 대상이지 정답 논리가 아니다.
  \item ③ 구름 부족이 기온 상승의 원인으로 제시되지 않았다. 지문의 인과 방향도 다르다.
  \item ④ 비 부족은 농업의 결과가 아니라 식물 생장을 어렵게 하는 환경 조건으로 제시된다.
  \item ⑤ 추위가 식물의 에너지 사용을 줄여 유리하다는 내용은 없다.
\end{itemize}

\bigskip

% ===== Q3 해설 =====
\qno{3}\textbf{[서술형 --- 표현 찾기]} \textbf{모범답안}

\vspace{6pt}
\begin{pbox}
\textbf{This didn't happen overnight.}
\end{pbox}

\vspace{6pt}
\noindent\textbf{채점 포인트}:
\begin{itemize}[leftmargin=2em,itemsep=2pt]
  \item \textbf{[2점]} \textit{This didn't happen overnight.}를 정확히 쓰면 정답
  \item \textbf{[1점]} 핵심 표현 \textit{didn't happen overnight}만 쓰면 부분 점수
\end{itemize}

\vspace{8pt}
\noindent{\sffamily\bfseries 4. 제목 - 정답: ①}
\par 지문은 빙하기의 춥고 건조한 환경에서는 농업이 어려웠지만, 기온 상승과 습윤화가 생명과 농업 가능성을 열었다는 흐름을 설명한다. 따라서 ①이 가장 적절하다.

\end{document}
